\chapter{Conclusions and Future Work}
\label{ch:conclusions}

This thesis set out to design and build a whistleblowing platform that protects a
source even against adversaries who can seize servers, coerce individuals, and
compel hosting providers --- and to do so without resting the source's safety on a
single trusted server or a single trusted chip. The preceding chapters described
the resulting system, \zkwns, from its theoretical foundations through its
architecture, implementation, and evaluation. This final chapter draws the threads
together, reflects on the central argument, and sets out a concrete path from
research prototype to production system.

\section{Summary of the Work}
\label{sec:concl-summary}

The work began from a clear problem: existing whistleblowing tools force a source
to trust something that a powerful adversary can attack --- a central server, an
organisation, or a piece of trusted hardware. Chapters~\ref{ch:background}
and~\ref{ch:related} grounded this in the relevant cryptography and in a
survey of prior systems, identifying a gap: no existing platform combined an
automatic, censorship-resistant trigger with key custody that relies on neither a
central party nor trusted hardware, while also offering credible-yet-anonymous
provenance and unlinkable rewards.

Chapter~\ref{ch:architecture} answered that gap with a four-layer design in which
the storage, key-custody, trigger, and client layers are independent decentralised
systems, and in which the trust root is deliberately moved from hardware to
mathematics. Chapter~\ref{ch:implementation} realised the design: in-browser
AES-256-GCM encryption, a minimal on-chain dead man's switch whose \isdeceased
predicate gates everything, threshold key custody via the Lit MPC network,
permanent storage on Arweave, ERC-5564 stealth-address payments, and zkTLS
provenance via Reclaim, all composed in a fixed, auditable pipeline that never lets
plaintext or an unsealed key leave the device. Chapter~\ref{ch:evaluation}
established functional correctness through 91 unit tests, mapped each threat to its
mitigation, and assessed the system honestly --- including the one property, live
validation of the release path, that remains to be demonstrated end to end.

\section{Contributions}
\label{sec:concl-contributions}

The thesis makes the following contributions.

\textbf{An integrated architecture for trust-minimised whistleblowing.} The
principal contribution is the design itself: a coherent composition of a dead
man's switch, threshold key custody, permanent storage, anonymous provenance, and
unlinkable payments into a single system with no central point of failure and no
reliance on trusted hardware. Each primitive existed before; assembling them so
that their guarantees reinforce rather than undermine one another is the
contribution.

\textbf{A concrete realisation of the hardware-to-mathematics trust shift.} The
work demonstrates, in running code, that key custody for a release-on-death system
can be founded on a threshold network rather than on a TEE, and articulates why
this is preferable for the threat model in question. The deliberate choice to pin
to the MPC path, rather than a hardware-backed alternative, makes the argument more
than rhetorical.

\textbf{A layered, configurable trust model.} Rather than pretending one mechanism
suits all deployments, the system makes its trust assumptions explicit and
selectable --- most visibly in the registry's three verification tiers, which range
from a fully trustless on-chain proof check to a development-only convenience, with
the production path being cryptographic. This makes the system's guarantees legible
to its users.

\textbf{A security-first reference implementation.} The accompanying code applies
defence-in-depth throughout: checks--effects--interactions with reentrancy guards
on every fund move, server-side confinement of secrets behind allow-lists and rate
limits, no-store handling of ciphertext, and a strict in-browser encryption
boundary. It serves as a worked example of building such a system carefully.

\section{Reflection on the Central Argument}
\label{sec:concl-reflection}

The recurring theme of this thesis is where trust ultimately lives. Every security
system has a trust root; the engineering question is whether that root is as small,
as distributed, and as hard to attack as the threat model demands. For a tool whose
users may face state-level adversaries, a root that is a single chip from a single
vendor --- one that has been broken before by precisely such adversaries --- is hard to
justify. Distributing the root across an independent quorum does not make trust
disappear, but it raises the bar from ``compromise one party'' to ``compromise a
threshold of independent parties simultaneously,'' which is the strongest position
available with deployable technology today.

That reframing --- from hiding a secret inside hardware to splitting it across a
network --- is the idea the system is built to embody. The prototype shows it is
buildable now, not merely conceivable.

\section{Future Work}
\label{sec:concl-future}

Several concrete directions would carry the prototype toward a production system.

\textbf{End-to-end validation of the release path.} The most immediate task is to
exercise the Lit~v8 decryption flow against a live Naga deployment on a supported
testnet such as Base Sepolia, confirming the full release-on-death cycle and
closing the one open item against design goal P4.

\textbf{Large-payload support.} Implementing the split-blob scheme --- a large
ciphertext stored separately with a compact, key-bearing manifest kept within the
free storage tier --- would lift the 100~KiB ceiling and support substantial leaks
such as document caches.

\textbf{Stronger metadata privacy.} Reducing the on-chain footprint of a vault ---
hiding the recipient, obscuring check-in timing, and decorrelating heartbeats from
external events --- would harden the system against traffic-analysis-style inference,
which the current design does not fully address.

\textbf{Liveness automation.} Heartbeats could be made less burdensome and less
error-prone through delegated or automated check-ins (for example, a trusted device
or a guardian quorum) that reduce the risk of an accidental trigger while
preserving the source's ultimate control.

\textbf{Hardening and assurance.} A production deployment would warrant a
professional security audit, broader fuzz and invariant testing, and ideally formal
verification of the dead man's switch and the marketplace settlement logic, given
their direct handling of trigger conditions and funds.

\textbf{Governance and resilience.} Finally, the platform's own parameters --- fees,
verifier addresses, supported providers --- and its long-term stewardship would
benefit from a transparent, decentralised governance model, so that the system that
removes single points of technical failure does not reintroduce a single point of
administrative failure.

\section{Closing Remarks}
\label{sec:concl-closing}

Whistleblowers expose wrongdoing at extraordinary personal risk, and the tools they
rely on should not ask them to trust a single party that a determined adversary can
break. \zkw is an argument, in design and in code, that they need not: that by
composing a censorship-resistant trigger, threshold key custody, anonymous
provenance, and unlinkable rewards, it is possible to build a system whose
protection survives the seizure of any one of its parts. The prototype does not
solve every problem --- its limitations are stated plainly --- but it demonstrates that
the harder, more trustworthy architecture is within reach with the cryptography
available today. Moving the trust root from a chip to a distributed mathematical
assumption is not merely a technical refinement; for the people these tools are
meant to protect, it is the difference between a promise that can be broken by one
actor and one that cannot.
